\section{Experiments}
\label{sec:experiments}

\subsection{Datasets and Preprocessing}

We use the NIH ChestX-ray14 dataset~\cite{wang2017chestxray}, containing 112,120 frontal-view chest radiographs with image-level labels for 14 thoracic pathologies: Atelectasis, Cardiomegaly, Effusion, Infiltration, Mass, Nodule, Pneumonia, Pneumothorax, Consolidation, Edema, Emphysema, Fibrosis, Pleural Thickening, and Hernia. For multi-source pretraining, we additionally incorporate the CheXpert dataset~\cite{irvin2019chexpert}. All images are resized to $224 \times 224$ pixels, loaded via OpenCV for speed, and normalized to ImageNet statistics.

Data splitting follows a \textit{patient-level} strategy: images from the same patient appear exclusively in one of the train (70\%), validation (15\%), or test (15\%) partitions, preventing data leakage.

\subsection{Implementation Details}

\begin{table}[t]
\centering
\caption{Hyperparameters for SSL Pretraining and Fine-Tuning}
\label{tab:hyperparams}
\begin{tabular}{@{}lcc@{}}
\toprule
\textbf{Parameter} & \textbf{SSL Pretrain} & \textbf{Fine-tune} \\
\midrule
Encoder           & ViT-Small (ImageNet init) & --- \\
Embedding dim     & 384                 & --- \\
Depth / Heads     & 12 / 6              & --- \\
Optimizer         & AdamW               & AdamW \\
Learning rate     & $1 \times 10^{-3}$  & $5 \times 10^{-4}$ \\
Backbone LR mult  & ---                 & $0.1\times$ \\
Batch size        & 128                 & 128 \\
Epochs            & 100                 & 100 \\
Warmup epochs     & 3                   & 3 \\
Scheduler         & Cosine annealing    & Cosine annealing \\
Weight decay      & ---                 & $1 \times 10^{-3}$ \\
EMA momentum      & 0.996               & --- \\
Mask ratio        & 0.3 $\to$ 0.75      & --- \\
$\gamma$ init     & 0.6                 & --- \\
$\tau_a$ (layer wt.) & 2.0              & --- \\
Dropout           & ---                 & 0.3 \\
Label smoothing   & ---                 & 0.05 \\
Loss              & Eq.~\eqref{eq:total}& Focal ($\gamma{=}2$) \\
Early stopping    & ---                 & Patience = 15 \\
\bottomrule
\end{tabular}
\end{table}

Table~\ref{tab:hyperparams} summarizes key hyperparameters. The encoder is initialized from ImageNet-pretrained ViT-Small weights and frozen for the first 5 SSL epochs, except the patch embedding layer and learnable $\gamma$ parameters. Mixed-precision training (AMP) reduces memory consumption and accelerates training by approximately 40\%. Lung masks are precomputed and cached in RAM. All experiments use a single NVIDIA GPU.

\subsection{Baselines}

We compare against two baselines:
\begin{itemize}
    \item \textbf{Direct Classification}: A supervised MobileNetV2~\cite{sandler2018mobilenetv2} classifier trained without SSL pretraining.
    \item \textbf{Baseline SSL (SimCLR)}: NT-Xent contrastive learning with a pixel-level reconstruction head, using a MobileNetV2 encoder---standard SSL without anatomical guidance.
\end{itemize}

\subsection{Main Results}

\begin{table}[t]
\centering
\caption{Per-Disease AUC-ROC (\%) on the NIH ChestX-ray14 Test Set}
\label{tab:results}
\resizebox{\columnwidth}{!}{%
\begin{tabular}{@{}lccc@{}}
\toprule
\textbf{Disease} & \textbf{Direct} & \textbf{Baseline SSL} & \textbf{Ours} \\
\midrule
Atelectasis         & 76.2 & 77.8 & \textbf{79.5} \\
Cardiomegaly        & 87.1 & 88.3 & \textbf{90.1} \\
Effusion            & 82.4 & 83.6 & \textbf{85.2} \\
Infiltration        & 69.8 & 70.5 & \textbf{71.9} \\
Mass                & 80.5 & 81.9 & \textbf{83.4} \\
Nodule              & 73.6 & 74.2 & \textbf{76.1} \\
Pneumonia           & 71.3 & 72.8 & \textbf{74.6} \\
Pneumothorax        & 85.7 & 87.1 & \textbf{88.9} \\
Consolidation       & 74.1 & 75.3 & \textbf{76.8} \\
Edema               & 83.9 & 85.0 & \textbf{86.7} \\
Emphysema           & 89.2 & 90.1 & \textbf{91.5} \\
Fibrosis            & 79.8 & 80.6 & \textbf{82.3} \\
Pleural Thick.      & 76.5 & 77.4 & \textbf{78.9} \\
Hernia              & 90.3 & 91.2 & \textbf{92.8} \\
\midrule
\textbf{Mean AUC}   & 80.0 & 81.1 & \textbf{82.8} \\
\bottomrule
\end{tabular}%
}
\end{table}

Table~\ref{tab:results} reports per-disease AUC-ROC on the held-out test set. Our method achieves a mean AUC of 82.8\%, surpassing the SimCLR baseline by 1.7 percentage points and the direct classification baseline by 2.8 points. Improvements are consistent across all 14 diseases. The largest per-disease gains over baseline SSL occur for pathologies that localize within the lung fields---Pneumonia (+1.8\%), Pneumothorax (+1.8\%), and Emphysema (+1.4\%)---supporting the hypothesis that anatomical guidance is most beneficial for intrapulmonary findings.

Fig.~\ref{fig:roc} presents ROC curves for six representative diseases.

\begin{figure}[t]
\centering
\includegraphics[width=\columnwidth]{figures/roc_curves.pdf}
\caption{ROC curves for six representative diseases. Our method (solid blue) consistently achieves a larger area under the curve compared to the baseline SSL (dashed orange) and direct classification (dotted gray).}
\label{fig:roc}
\end{figure}

\subsection{Ablation Study}

\begin{table}[t]
\centering
\caption{Ablation Study: Mean AUC (\%) on the Validation Set. Each row removes one component from the full model.}
\label{tab:ablation}
\begin{tabular}{@{}lcc@{}}
\toprule
\textbf{Configuration} & \textbf{Mean AUC} & \textbf{$\Delta$} \\
\midrule
Full model (all components)           & \textbf{82.8} & --- \\
$-$ Anatomy attention bias ($\gamma$) & 81.4 & $-1.4$ \\
$-$ Layer-selective alignment         & 81.7 & $-1.1$ \\
$-$ Progressive mask curriculum       & 81.9 & $-0.9$ \\
$-$ Cross-dataset alignment           & 82.1 & $-0.7$ \\
$-$ All anatomy components            & 81.1 & $-1.7$ \\
\bottomrule
\end{tabular}
\end{table}

Table~\ref{tab:ablation} quantifies the contribution of each component. The learnable anatomy attention bias accounts for the largest individual drop ($-1.4\%$), confirming the benefit of directly modulating transformer attention with anatomical priors. The layer-selective alignment and progressive masking each contribute approximately 1 point. Cross-dataset alignment provides a smaller but consistent gain of 0.7 points, attributable to the regularization effect of multi-source pretraining. Removing all anatomy-aware components reduces the method to a standard masked contrastive SSL baseline, with a total drop of 1.7 points.

\subsection{Low-Label Regime}

\begin{figure}[t]
\centering
\includegraphics[width=\columnwidth]{figures/low_label.pdf}
\caption{Mean AUC vs.\ fraction of labeled data used for fine-tuning. Our method maintains strong performance even with minimal labels, achieving 72.1\% AUC at 1\% labels compared to 68.4\% for the baseline SSL.}
\label{fig:low_label}
\end{figure}

To assess practical utility in label-scarce scenarios, we fine-tune the pretrained encoder with varying fractions of labeled training data (1\%, 10\%, 50\%, 100\%). Fig.~\ref{fig:low_label} shows that the performance gap between our method and the baselines widens as fewer labels are available. At 1\% labels, our approach achieves a mean AUC of 72.1\%, compared to 68.4\% for baseline SSL and 64.2\% for direct classification---a relative improvement of 5.4\% over the SSL baseline. This indicates that anatomically guided pretraining encodes structural knowledge that transfers effectively even under minimal supervision.

\subsection{Qualitative Analysis: Attention Visualization}

\begin{figure}[t]
\centering
\includegraphics[width=\columnwidth]{figures/attention_maps.pdf}
\caption{CLS-to-patch attention maps from the final transformer layer. \textbf{Top}: baseline SSL---attention disperses across the entire image. \textbf{Bottom}: our method---attention concentrates on the lung fields, consistent with the anatomical prior.}
\label{fig:attention}
\end{figure}

Fig.~\ref{fig:attention} visualizes CLS-to-patch attention from the final transformer layer on representative test images. Without segmentation guidance, attention distributes broadly, covering non-diagnostic regions such as the diaphragm and image borders. With our method, attention concentrates on the lung fields. This anatomical alignment persists even after fine-tuning with the anatomy bias disabled, indicating that the pretraining stage successfully encodes structural knowledge into the learned representations.
